%! Trigonometric Equations and Inequalities — Unit 7, Lesson 7.3 — Warm-Up (Answer Key)
\documentclass[10pt]{article}
\usepackage{precalculus-article}
\usepackage{precalculus-key}   % pulls in -boxes; do NOT also load -boxes
\newcommand{\TallMath}[1]{$\displaystyle #1\rule[-1.4em]{0pt}{3.2em}$}

\begin{document}

\pageheader{Unit 7, Lesson 7.3}{Warm-Up --- Answer Key}
\namedateperiod

\vspace{0.10in}

\begin{notesbox}{Spiral Review --- Prerequisites for Lesson 7.3}

\textbf{1.}\quad Evaluate each inverse trigonometric expression.
Give your answer in radians.

\vspace{4pt}
(a)\quad $\arcsin\!\left(\dfrac{\sqrt{2}}{2}\right) = $ \ans{$\dfrac{\pi}{4}$}
\hspace{1cm}
(b)\quad $\arccos\!\left(-\dfrac{1}{2}\right) = $ \ans{$\dfrac{2\pi}{3}$}

\vspace{0.14in}
\textbf{2.}\quad Given that $\sin\theta = \dfrac{1}{2}$, list \textbf{all}
values of $\theta$ in the interval $[0, 2\pi)$ that satisfy this equation.

\vspace{4pt}
$\theta = $ \ans{$\dfrac{\pi}{6}$ and $\dfrac{5\pi}{6}$}

\vspace{4pt}
Explain briefly why there are two solutions (not just one):
\ans{Sine is positive in both Quadrant I and Quadrant II. The reference angle $\tfrac{\pi}{6}$ gives one solution in Q1; the supplementary angle $\pi - \tfrac{\pi}{6} = \tfrac{5\pi}{6}$ gives the Q2 solution.}

\vspace{0.14in}
\textbf{3.}\quad Solve for $\sin x$ by isolating the trig function.
\textbf{Do not} find the actual angle --- just isolate $\sin x$.

\vspace{4pt}
\hspace{1em}$2\sin x + 1 = 0$

\vspace{6pt}
$\sin x = $ \ans{$-\dfrac{1}{2}$}

\vspace{0.14in}
\textbf{4.}\quad Which of the following expressions gives the value of
$\theta$ in Quadrant II if $\sin\theta = \dfrac{3}{5}$?
Circle the correct answer.

\vspace{4pt}
\hspace{2em}
$\displaystyle\theta = \arcsin\!\!\left(\tfrac{3}{5}\right)$
\hspace{2em}
\ans{$\displaystyle\theta = \pi - \arcsin\!\!\left(\tfrac{3}{5}\right)$}
\hspace{2em}
$\displaystyle\theta = \pi + \arcsin\!\!\left(\tfrac{3}{5}\right)$

\begin{teachernote}
Q4 note: $\arcsin(3/5)$ gives the Q1 reference angle. Q2 reflection is
$\pi - \arcsin(3/5)$. The third option $\pi + \arcsin(3/5)$ lands in Q3
where sine is negative.
\end{teachernote}

\end{notesbox}

\end{document}
